\begin{abstract}
We prove Brennan's conjecture: if a simply connected plane domain admits a
conformal bijection $\varphi$ onto the unit disk, then $|\varphi'|^s$ is
area-integrable for every $4/3<s<4$. We also prove the sharp inverse-square
integral-means exponent $B_{\SSS}(-2)=1$ for the normalized schlicht class
$\SSS$.
\end{abstract}

\section{Introduction}
\label{sec:introduction}

Let $\DD=\{z\in\CC:|z|<1\}$, and write $dA=dx\,dy$ for planar
Lebesgue measure. A domain is a connected open subset of $\CC$; a
univalent map is an injective holomorphic map. Brennan's conjecture asks
whether, for every conformal bijection $\varphi:W\to\DD$ from a simply
connected plane domain, the derivative satisfies
\[
 \int_W|\varphi'(z)|^s\dd A(z)<\infty
 \qquad\left(\frac43<s<4\right).
\]
There is no regularity assumption on the boundary and no boundedness
assumption on $W$. The problem measures how strongly a conformal map can
compress area near an arbitrary boundary. The conjecture dates to
Brennan's 1978 work \cite{Brennan1978}; the domain and integral-means
formulations are recorded in \cite[Chapter~3]{Bertilsson1999}.

The inverse map $f=\varphi^{-1}:\DD\to W$ puts all source domains on
the same footing. Conformal change of variables gives
\begin{equation}\label{eq:inverse-area}
 \int_W|\varphi'|^s\dd A
 =\int_{\DD}|f'|^{2-s}\dd A.
\end{equation}
Thus the conjectured interval becomes $-2<t<2/3$, with $t=2-s$.
Positive powers detect large derivatives, whereas the difficult negative
powers detect small derivatives. The inverse-square power $t=-2$ is
the limiting case relevant to the upper endpoint $s=4$.

To state its radial form, let
\[
 \SSS=\{f:\DD\to\CC\text{ univalent}:f(0)=0,\ f'(0)=1\}
\]
be the normalized schlicht class. For $f\in\SSS$, $t\in\mathbb R$,
and $0<r<1$, define
\[
 M_t[f'](r)=\frac1{2\pi}\int_{-\pi}^{\pi}
                  |f'(re^{i\theta})|^t\dd\theta,
\]
and put
\begin{align*}
 \beta_f(t)&=\limsup_{r\uparrow1}
       \frac{\log M_t[f'](r)}{\log(1/(1-r))},\\
 B_{\SSS}(t)&=\sup_{f\in\SSS}\beta_f(t).
\end{align*}
Equivalently, $\beta_f(t)$ is the infimum of the real $b$ for which
$M_t[f'](r)=O((1-r)^{-b})$ as $r\uparrow1$; the implicit constant may
depend on $f$. In particular, the supremum defining $B_{\SSS}$ is taken after forming
each map's limiting exponent. We prove the following uniform estimate
directly, and then derive the area assertion.

\begin{theorem}[Brennan's conjecture and inverse-square means]
\label{thm:main}
For every $\varepsilon>0$ there is $C_\varepsilon<\infty$ such that
\begin{equation}\label{eq:main-uniform}
 M_{-2}[f'](r)\le C_\varepsilon(1-r)^{-1-\varepsilon}
 \qquad\left(f\in\SSS,\ \frac12\le r<1\right).
\end{equation}
Moreover, $B_{\SSS}(-2)=1$.

Let $W\subset\CC$ be any simply connected domain whose boundary in
the Riemann sphere contains at least two points. Every conformal
bijection $\varphi:W\to\DD$ satisfies
\[
 \int_W|\varphi'(z)|^s\dd A(z)<\infty
 \qquad\left(\frac43<s<4\right).
\]
Equivalently, every univalent map $f:\DD\to\CC$ satisfies
\[
 \int_{\DD}|f'(w)|^t\dd A(w)<\infty
 \qquad\left(-2<t<\frac23\right).
\]
\end{theorem}

The domain correspondence, including the boundary condition, is
justified in Section~\ref{sec:integrability}. That section also derives
$B_{\SSS}(t)=|t|-1$ for every $t\le-2$; the corresponding uniform
radial bounds retain an arbitrary positive loss. The radial exponent $1$
is compatible with divergence of the area integral at $t=-2$.
The other area endpoint also fails for the same classical example.

\begin{corollary}[Sharpness of the open range]\label{cor:sharp-endpoints}
The Koebe function $f(z)=z/(1-z)^2$ belongs to $\SSS$ and satisfies
\[
 \int_{\DD}|f'|^{-2}\dd A=\infty,
 \qquad
 \int_{\DD}|f'|^{2/3}\dd A=\infty.
\]
Its inverse therefore fails area integrability at both $s=4$ and $s=4/3$.
\end{corollary}

